\documentclass[aip,reprint,groupedaddress,onecolumn,superscriptaddress]{revtex4-2}
\usepackage{graphicx}
\usepackage{dcolumn}
\usepackage{bm}
\usepackage{amsmath}
\usepackage{siunitx}
\usepackage{booktabs}

\begin{document}

\title{Bulletpoint List for Conclusion}
\date{\today}

\maketitle

\section{Summary: A Proof-of-Principle Demonstration}
\begin{itemize}
    \item We have built and operated a prototype demonstrator that creates a direct metrological link between a quantum standard (a single trapped $^{138}$Ba$^+$ ion) and the ambient refractive index at 1762 nm.
    \item The key output is a set of modified Edlén coefficients for 1762 nm, with an uncertainty budget derived from and traceable to the quantum standard itself.
    \item This work provides the first experimental validation that a qubit's own frequency reference can be used to characterize the environmental parameters that affect it, addressing the initial step of the "traceability gap."
\end{itemize}

\section{Demonstrated Impact \& Paradigm Potential}
\begin{itemize}
    \item \textbf{For Quantum Information/Control:} The prototype provides a concrete, co-design pathway to move beyond indirect atmospheric models. It directly tackles a known bottleneck (phase instability) for the fidelity and scalability of field-deployable quantum networks and sensors.
    \item \textbf{Conceptual Shift:} It demonstrates a viable framework for a paradigm shift—from treating the environment as mere noise to suppress, toward precisely characterizing it as a quantum-traceable quantity. This changes how we define performance limits for real-world quantum systems.
\end{itemize}

\section{Clear Outlook \& Next Steps}
\begin{itemize}
    \item \textbf{Immediate Technical Path}: The co-design principle is directly adaptable to other qubit platforms (e.g., neutral atoms, solid-state defects) to establish traceability for their specific operational wavelengths.
    \item \textbf{Development Goal}: The next critical phase is the development of compact, automated sensor nodes for practical integration into functional quantum network architectures.
    \item \textbf{Broader Metrological Impact}: This methodology opens a path for quantum metrology to expand its domain beyond time/frequency, potentially contributing to fundamental atmospheric science by providing benchmark data for model validation.
\end{itemize}

\newpage

\section*{Conclusion}

In summary, we have demonstrated a proof-of-principle for creating a direct metrological link between a primary quantum frequency standard and a key environmental parameter. Using a single trapped $^{138}$Ba$^+$ ion as a reference, we implemented a co-designed apparatus to provide SI traceability for a dual-wavelength interferometer. This prototype enabled a measurement of the refractive index of air at 1762 nm, yielding modified Edlén coefficients with an uncertainty budget derived from the quantum standard. This result constitutes a first step in bridging the traceability gap, showing that parameters critical for quantum network operation can be validated by the qubit platform itself.

The significance of this demonstration lies in establishing a viable framework for endowing quantum systems with intrinsic environmental traceability. It suggests a paradigm shift from merely suppressing environmental noise as an uncontrolled disturbance, toward precisely characterizing it as a measurable quantity using the system's own quantum resources. For the quantum information and control community, this prototype provides a concrete pathway to move beyond reliance on indirect atmospheric models, directly addressing a known bottleneck for fidelity and scalability in field-deployable systems.

Looking forward, the co-design principle demonstrated here can be adapted to other leading qubit platforms, such as neutral atoms or solid-state defects, to establish traceability for their respective operational wavelengths. Further development towards compact, automated sensor nodes will be essential for integration into functional quantum networks. Moreover, this methodology opens a path for quantum metrology to expand its domain, potentially contributing to fundamental atmospheric science by providing benchmark data for model validation at specific wavelengths. The realization of such quantum-traceable environmental sensing represents a critical enabling step for scalable quantum technologies that must operate outside controlled laboratory conditions.

\bibliography{reference}

\end{document}